\begin{lemma}[The final output of a positive series is an ETT]
Let \(G\) be a finite loopless multigraph, let \(e_0\) be an edge, let
\(\varphi\) be a partial \(k\)-edge-coloring, and let \(\mathcal S\) be a
certified Tashkinov series from \((G,e_0,\varphi)\) with \(r>0\)
noninitial iterations.  Write \(T_{r+1}\) for the final tree stored by
\(\mathcal S\) and \(\varphi_r\) for its terminal coloring.  Then there is
an extended Tashkinov tree displayed on exactly \(T_{r+1}\), with generating
coloring exactly \(\varphi_r\), whose rung is exactly \(r\) and whose
dependent certified series is the given series \(\mathcal S\).
\end{lemma}
\begin{proof}
By the definition of a certified series in \noderef{TashkinovSeries}, the
series has a \noderef{TashkinovSeriesStep} certificate at every index from
zero through \(r\).  Fix \(i<r\) and inspect the certificate at the positive
index \(i+1\).  It cannot be the initial constructor: that constructor
certifies index zero, whereas \(i+1\) is positive.  The remaining three
constructors are SE, PE, and RE.  In each of them, the two final fields state
respectively that \(T_{i+1}\) is an ordered initial segment of
\(T_{i+2}\) and that
\[
                         |T_{i+1}|<|T_{i+2}|.
\]
Thus both assertions hold for every \(i<r\).

Since \(r>0\), take \(i=r-1\).  The preceding conclusion then gives the
ordered initial-segment relation
\(T_r\subseteq T_{r+1}\) together with the strict length inequality
\(|T_r|<|T_{r+1}|\).  Hence this containment is strict.  Moreover,
\(T_{r+1}\) is an ordered initial segment of itself: in the definition
\noderef{TreeSequenceInitialSegment}, choose the reflexive length inequality,
and all required vertex and edge equalities are reflexive.

Now apply the definition \noderef{ExtendedTashkinovTree} with displayed tree
\(T_{r+1}\), rung \(r\), and the original series \(\mathcal S\).  The strict
lower containment is the relation and inequality just obtained, the upper
containment is the reflexive certificate just constructed, and the displayed
generating coloring is \(\varphi_r\) by reflexivity.  This constructs the
required ETT.  Because its rung and its dependent series were assigned
directly to \(r\) and \(\mathcal S\), the asserted rung equality and dependent
series agreement are both reflexive.
\end{proof}
